\section{The Instrument That Lied}
\label{sec:instrument-lied}

Everything in the previous section was, at one point, wrong in the same way: the measurements
were taken on a time axis that did not exist. This section is the most important one in the
chapter, because the failure was silent, it survived a code review, and it certified a defect
as fixed.

\subsection{The Sampler That Claimed 400\,Hz and Ran at 158}

The loop was written to sample at 400\,Hz and handed the detector a constant:

\begin{minted}{rust}
const SAMPLE_PERIOD_CYCLES: u32 = 47_500;  // 400 Hz at 19 MHz
const DT_US: u32 = 2_500;                  // 400 Hz, in microseconds
...
application.sample_us(motion.accel_mg, DT_US);
\end{minted}

The loop was paced correctly — it never sampled \emph{faster} than 400\,Hz. But each iteration
also did work: a 12-byte I2C burst, a detector update, a record write, and periodically a
1090-byte display flush over the same bit-banged bus. The true interval was set by the
\textbf{slowest} part of the body, not by the period:

\begin{minted}{text}
measured: 317 samples in 2 s  ->  158 Hz, mean interval 6.3 ms
\end{minted}

Every timeout in the detector was therefore 2.5 times too short, and every duration in the
first analysis — ``40 ms of free fall'' — was 2.5 times too small. The measurements were
self-consistent, which is exactly why nobody noticed.

\subsection{What Gave It Away}

Not a test. A pattern in the output:

\begin{minted}{text}
AIRBORNE  free 85ms   ->  FALL   after 0.62 s
AIRBORNE  free 52ms   ->  FALL   after 0.62 s
AIRBORNE  free 120ms  ->  FALL   after 0.62 s
AIRBORNE  free 40ms   ->  FALL   after 0.53 s
\end{minted}

Four different drops, four different free falls, and the settle time is the same to the
millisecond. \textbf{The number was describing the code, not the world.} That is the signature
of a broken instrument: when an output is constant across inputs that should vary, you are
reading your own parameter back.

\subsection{Three Fixes, and One of Them Was Deleted by the Compiler}

\begin{enumerate}
    \item \textbf{Measure the interval.} The loop takes the difference of the DWT cycle
          counter between samples and passes that everywhere — detector, recorder, alarm.
          \begin{minted}{rust}
let dt_us = now.wrapping_sub(last_sample_cycles) / CYCLES_PER_US;
last_sample_cycles = now;
application.sample_us(motion.accel_mg, dt_us);
          \end{minted}
    \item \textbf{Record the interval with every sample.} A dataset whose timing is unknown
          can only be analysed dishonestly. One byte per sample, in 250\,$\mu$s units.
    \item \textbf{Get the display out of the critical path.} Refresh every 800 samples instead
          of 80, and never while the device is airborne or within 1.2\,s of an event: a flush
          in the middle of a fall stalls the sampler exactly when the samples matter. Measured
          effect: 158\,Hz to 235\,Hz without touching the sensor.
\end{enumerate}

And then the second-order failure, which is the reason this section exists:

\begin{warningbox}
    \textbf{The compiler deleted the per-sample timing array.} The recording was written to
    ``carry its own timing'', and nothing in the crate ever read that array — so the optimiser
    removed it from the image entirely. The code compiled, the comment claimed a property, the
    symbol was absent from the binary, and the recording carried nothing. It was found by
    running \codeinline{arm-none-eabi-nm} and looking for the symbol instead of trusting the
    comment. The writes are volatile now, and the check is part of the habit: \emph{a property
    you cannot see in the artefact is a property you do not have.}
\end{warningbox}

\subsection{The Instrument Rules}

\begin{keytakeaways}
    \item \textbf{Never pass the interval you asked for; pass the one you measured.} The
          intended rate is a wish. The DWT counter is the fact.
    \item \textbf{A constant output across varying inputs means you are measuring yourself.}
          The 0.62\,s settle was a parameter echoed back.
    \item \textbf{A recording must carry its timing.} Otherwise every later analysis inherits
          an assumption nobody remembers making.
    \item \textbf{Verify the artefact, not the intent.} Check that the symbol is in the binary,
          that the counter advanced, that the file grew. Comments describe intentions; the
          binary is what runs.
    \item \textbf{Calibrate the instrument before calibrating the product.} The rate, the
          units and the resolution have to be right before any threshold on top of them means
          anything.
\end{keytakeaways}
